\CNsection{}{معرفی}

\textbf{شبکه‌های ارتباطی پیشرفته} یک درس ۳ واحدی (۴۸ ساعت، ۱۶ هفته \ensuremath{\times} ۳ ساعت) برای دانشجویان سال آخر کارشناسی و تحصیلات تکمیلی مهندسی برق و مخابرات است که به‌صورت سخنرانی همراه با آزمایشگاه و پروژه ارایه می‌شود. این درس یک سفر فنی و یکپارچه در تحول شبکه‌های موبایل از \lr{\mbox{GSM (2G)}} تا \lr{\mbox{LTE (4G)}}، \lr{5G NR/\allowbreak{}SA} و مفاهیم نوظهور \lr{\mbox{6G}} است \lr{—} نه مجموعه‌ای از فناوری‌های جدا، بلکه تحلیل پیوسته‌ای از این‌که \textbf{چرا} هر نسل پدید آمد، \textbf{چه} مشکلی را حل کرد و \textbf{چگونه} ترافیک از معماری آن عبور می‌کند. درس بر منطق «چرا /\allowbreak{} چه چیزی /\allowbreak{} چگونه» استوار است تا دانشجو به‌جای حفظ فهرست پروتکل‌ها، بتواند تصمیم‌های معماری را تحلیل و ارزیابی کند.

\CNsection{}{ساختار و محتوا}

مطالب در \textbf{۲۶ ماژول} سازمان یافته است:

\begin{itemize}
\tightlist
\item
  \textbf{پایه ۲\lr{\mbox{G}}/\allowbreak{}۳\lr{\mbox{G}} (ماژول‌های ۱\lr{–}۴):} معماری \lr{\mbox{GSM}}، \lr{\mbox{GPRS}} و \lr{\mbox{EDGE}}، سوییچینگ مداری و سیگنالینگ \lr{\mbox{SS7}}، شبکه‌های هسته \lr{2G/\allowbreak{}3G} و مبانی مدیریت تحرک و تحویل \lr{\mbox{(handover)}}.
\item
  \textbf{دسترسی و هسته \lr{\mbox{4G}} (ماژول‌های ۵\lr{–}۱۳):} معماری \lr{\mbox{LTE}}، شبکه دسترسی \lr{\mbox{E-UTRAN}}، پشته پروتکل، زمان‌بندی منابع، کانال دسترسی تصادفی \lr{\mbox{(RACH)}}، هسته بسته تکامل‌یافته \lr{\mbox{(EPC)}}، پروتکل \lr{\mbox{Diameter}}، \lr{\mbox{HSS}} و احراز هویت، \lr{\mbox{MME}}، و \lr{S-GW/\allowbreak{}P-GW} با تونل‌سازی \lr{\mbox{GTP}}.
\item
  \textbf{معماری‌های نرم‌افزارمحور (ماژول‌های ۱۵\lr{–}۱۸):} شبکه‌سازی نرم‌افزارمحور \lr{\mbox{(SDN)}}، مجازی‌سازی توابع شبکه \lr{\mbox{(NFV)}}، فرایند استانداردسازی \lr{\mbox{3GPP}} و معماری خدمات‌محور \lr{\mbox{5G (SBA)}}.
\item
  \textbf{هسته و رویه‌های \lr{\mbox{5G}} (ماژول‌های ۱۴، ۱۹\lr{–}۲۲):} معماری هسته \lr{\mbox{5G}}، رویه‌های سیگنالینگ (ثبت‌نام، احراز هویت \lr{\mbox{5G-AKA}}، برقراری و آزادسازی نشست \lr{\mbox{PDU}}، صفحه‌بندی، تحویل)، شبکه‌بندی و تحول چارچوب \lr{\mbox{QoS}} (از \lr{\mbox{QCI}} به \lr{\mbox{5QI}}).
\item
  \textbf{صفحه‌ها و چشم‌انداز (ماژول‌های ۲۳\lr{–}۲۶):} صفحه کاربر و صفحه کنترل، ترکیب جامع تحول شبکه، و جهت‌های پژوهشی \lr{\mbox{6G}} شامل \lr{\mbox{ISAC}}، \lr{\mbox{RIS}}، شبکه‌های بومی-هوش مصنوعی، \lr{\mbox{MIMO}} بدون سلول و \lr{\mbox{sub-THz}}.
\end{itemize}

هر ماژول با جدول‌های مقایسه‌ای، نمودارهای معماری، نمودارهای توالی سیگنالینگ گام‌به‌گام و پرسش‌های مرور همراه است. یک سند مرور پیش‌نیازها نیز مبانی \lr{\mbox{OSI}}، \lr{\mbox{RF}}، مدولاسیون، \lr{\mbox{OFDM}}، شبکه‌سازی \lr{\mbox{IP}}، احتمال و ریاضیات دسی‌بل را یادآوری می‌کند.

\begin{CNbox}[unbreakable]{obj}{اهداف یادگیری}

در پایان درس، دانشجو می‌تواند سیر تحول نسل‌ها را ترسیم کند؛ معماری‌های هسته (سوییچینگ مداری، بسته‌ای، \lr{\mbox{EPC}} تمام-\lr{\mbox{IP}} و هسته خدمات‌محور \lr{\mbox{5G}}) را تحلیل کند؛ مدیریت تحرک و تحویل را در نسل‌ها مقایسه کند؛ شبکه‌بندی، \lr{\mbox{QoS}}، \lr{\mbox{SDN}} و \lr{\mbox{NFV}} را در شبکه‌های مدرن به کار گیرد؛ رویه‌های هسته \lr{\mbox{5G}} را در سطح سیگنالینگ توضیح دهد؛ و جهت‌های پژوهشی \lr{\mbox{6G}} را به‌صورت انتقادی ارزیابی کند.

\end{CNbox}

\Needspace{15\baselineskip}

\CNsection{}{ارزشیابی و کار عملی}

{\def\LTcaptype{none} % do not increment counter
\Needspace{9\baselineskip}
\begin{longtable}[c]{@{}cc@{}}
\toprule\noalign{}
\rowcolor{CNink}\CNth{مؤلفه} & \CNth{وزن} \\
\CNheadrulesep
\endhead
\bottomrule\noalign{}
\endlastfoot
آزمون‌ها (میان‌ترم ۲۰٪ + پایان‌ترم ۳۰٪) & ۵۰٪ \\
آزمایشگاه‌ها (۶ تمرین \ensuremath{\times} ۵٪) & ۳۰٪ \\
پروژه (طرح پیشنهادی ۵٪ + گزارش و دمو ۱۵٪) & ۲۰٪ \\
\end{longtable}
}

تمرین‌های آزمایشگاهی شامل بودجه لینک \lr{\mbox{LTE}}، ردیابی سیگنالینگ \lr{\mbox{EPC}}، زمان‌بندی و تخصیص منابع \lr{\mbox{5G NR}}، کشف سرویس \lr{\mbox{SBA}}، پیکربندی شبکه‌بندی و ردیابی جریان ترافیک صفحه کاربر است. پروژه به‌صورت انفرادی یا دو نفره روی موضوعاتی مانند تحلیل مقایسه‌ای \lr{\mbox{LTE}} در برابر \lr{\mbox{5G}}، طراحی شبکه‌بندی سازمانی، پیاده‌سازی الگوریتم زمان‌بندی با شبیه‌ساز، یا تحلیل سیگنالینگ با فایل‌های ضبط بسته واقعی تعریف می‌شود.

\CNsection{}{منابع اصلی}

\lr{\mbox{Sesia}} و همکاران (\lr{\emph{\lr{LTE – The UMTS Long Term Evolution}}})، \lr{\mbox{Dahlman}} و همکاران (\lr{\emph{\lr{5G NR: The Next Generation Wireless Access Technology}}})، \lr{\mbox{Holma}} و \lr{\mbox{Toskala}} (\lr{\emph{\lr{LTE for UMTS}}})، مشخصات \lr{\mbox{3GPP}} (\lr{TS 23.501}، \CNnum{23.502}، \CNnum{24.501}، \CNnum{38.300)} و \lr{Rappaport (\emph{\lr{Wireless Communications}})}.
